\documentclass{article}

% Language setting
% Replace `english' with e.g. `spanish' to change the document language
\usepackage[english]{babel}

% Set page size and margins
% Replace `letterpaper' with `a4paper' for UK/EU standard size
\usepackage[letterpaper,top=2cm,bottom=2cm,left=3cm,right=3cm,marginparwidth=1.75cm]{geometry}

% Useful packages
\usepackage{amsmath}
\usepackage{amsthm}
\usepackage{amssymb}
\usepackage{graphicx}
\usepackage{algpseudocode}
\usepackage{algorithm}
\usepackage[colorlinks=true, allcolors=blue]{hyperref}
\renewcommand{\theenumi}{(\Alph{enumi})}

\usepackage{graphicx}
\graphicspath{ {./graph/} }


\newtheorem{proposition}{Proposition}[section]
\newtheorem{assumption}{Assumption}[section]
\newtheorem{theorem}{Theorem}
\newtheorem{definition}{Definition}

\title{Research note}

\begin{document}
\maketitle

\begin{abstract}
This is a research note for implementing quasi-Newton method on minimax problem.
\end{abstract}

\section{Introduction}
We focus on the minimax problem:
\begin{equation}
    \min_{x \in \mathbb{R}^n} \max_{y \in \mathbb{R}^m} f(x,y) = \min_{x \in \mathbb{R}^n} \Phi(x)
\end{equation}
where $f: \mathbb{R}^{n+m} \to \mathbb{R}$ is twice continuously differentiable, i.e. $f \in \mathcal{C}^2$.
We define the local optimal solution of this problem as follow:

\begin{definition}[strict local minimax(SLmM), Evtushenko (1974a)] $(x^*, y^*)$ is a strict local minimax point (SLmM) of a twice differentiable function $f$ if it is a stationary point and $\partial_{yy}f(x^*,y^*) < 0$ and $\mathbf{D}_{xx}f(x^*,y^*) > 0$
    
\end{definition}

\section{Preliminary}
We focus on the minimax problem:
\begin{equation}
    \min_{x \in \mathbb{R}^n} \max_{y \in \mathbb{R}^m} f(x,y) = \min_{x \in \mathbb{R}^n} \Phi(x)
\end{equation}
where $\Phi(x) = \max_{y \in \mathbb{R}^m} f(x,y)$. 
Suppose we can find a solution of the maximization problem $y(x) = \arg\max_{y \in \mathbb{R}^m} f(x,y)$, we have $\Phi(x) = f(x, y(x))$, and we can see the minimax problem as a minimization problem $\min_{x \in \mathbb{R}^n} \Phi(x)$ and a maximization problem $\max_{y \in \mathbb{R}^m} f(x,y)$.
\\
For the maximization problem, we have many ways to find the approximate solution such as gradient ascent, Newton method and quasi-Newton method.
\\
For the minimization problem, we need to analyze the function $\Phi(x) = f(x, y(x))$ furthermore.
We know that 
\begin{equation}
    \nabla \Phi(x) = \nabla_xf(x, y(x)) + \nabla y(x) \nabla_yf(x, y(x)) = \nabla_xf(x, y(x))
\end{equation}
since $y(x)$ solve $\max_{y \in \mathbb{R}^m}f(x,y)$ so $\nabla_yf(x,y(x)) = 0$.
Therefore, for the gradient descent step, the update will be $x^{k+1} = x^k - \alpha_k \nabla\Phi(x^k) = x^k - \alpha_k\nabla_xf(x^k, y(x^k))$.
If we want to take a Newton update, we need to calculate the Hessian $\nabla^2\Phi(x)$. From $\nabla \Phi(x) = \nabla_xf(x, y(x))$, we know that 
\begin{flalign}
    \label{eq3}
    \nabla^2\Phi(x) &= \nabla_{xx}f(x, y(x)) + \nabla y(x) \nabla_{yx} f(x, y(x)) \nonumber\\ 
                    &= \nabla_{xx}f(x, y(x)) - \nabla_{xy}f(x, y(x))(\nabla_{yy}f(x, y(x)))^{-1}\nabla_{yx}f(x, y(x))
\end{flalign}
The second equality use the fact that 
\begin{flalign*}
    &\nabla_{xy} f(x, y(x)) + \nabla y(x) \nabla_{yy}f(x,y(x)) = 0 \\
    &\nabla y(x) = -\nabla_{xy}f(x, y(x)) (\nabla_{yy}f(x, y(x)))^{-1}
\end{flalign*}
since $\nabla_yf(x, y(x)) = 0$. By using \hyperref[eq3]{equation 3}, the Newton update will be:
\begin{equation}
    x^{k+1} = x^k - (\nabla^2 \Phi(x^k))^{-1}\nabla\Phi(x^k)
\end{equation}
Since the Hessian used in Newton method might be time consuming, we can try to use quasi-Newton method to approximate the Hessian we need:
\begin{equation}
    x^{k+1} = x^k - H(x^k)^{-1}\nabla\Phi(x^k) = x^k - B(x^k)\nabla\Phi(x^k)
\end{equation}
where $H(x)$ is the approximate Hessian and $B(x)$ is the inverse of approximate Hessian.
Here, we have many option for the approximation such as BFGS.

\section{BFGS}

In BFGS, we approximate the Hessian by the step difference and the gradient difference.
Normal BFGS update for minimization problem $\min_{x \in \mathbb{R}^n}f(x)$ is 
\begin{equation}
    x_{k+1} = x_k - H_k \nabla f_k
\end{equation}
where $H$ is the inverse Hessian approximation:
\begin{equation}
    H_{k+1} = (I - \frac{s_k y_k^T}{y_k^T s_k})H_k(I - \frac{y_k s_k^T}{y_k^T s_k}) + \frac{s_k s_k^T}{y_k^T s_k},\; s_k = x_{k+1} - x_k,\; y_k = \nabla f_{k+1} - \nabla f_k
\end{equation}

So the total update for minimax problem is:
\begin{algorithm}
    \caption{BFGS for minimization problem}
    \begin{algorithmic}
        \State step size $\alpha$
        \State $P_{x,k} = -H_{x,k} \nabla_x f(x_k, y_k)$
        \State $s_{x,k} = \alpha P_{x,k}$
        \State $x_{k+1} = x_k + s_{x,k}$
        \State $g_{x,k} = \nabla_x f(x_{k+1}, y_k) - \nabla_x f(x_k, y_k)$
        \State $H_{x,k+1} = (I - \frac{s_{x,k} g_{x,k}^T}{g_{x,k}^T s_{x,k}})H_{x,k}(I - \frac{g_{x,k} s_{x,k}^T}{g_{x,k}^T s_{x,k}}) + \frac{s_{x,k} s_{x,k}^T}{g_{x,k}^T s_{x,k}}$
    \end{algorithmic}
\end{algorithm}
\begin{algorithm}
    \caption{BFGS update for maximization problem}
    \begin{algorithmic}
        \State step size $\alpha$
        \State $P_{y,k} = H_{y,k} \nabla_y f(x_k+1, y_k)$
        \State $s_{y,k} = \alpha P_{y,k}$
        \State $y_{k+1} = y_k + s_{y,k}$
        \State $g_{y,k} = \nabla_y f(x_{k+1}, y_{k+1}) - \nabla_y f(x_{k+1}, y_k)$
        \State $H_{y,k+1} = (I - \frac{s_{y,k} g_{y,k}^T}{g_{y,k}^T s_{y,k}})H_{y,k}(I - \frac{g_{y,k} s_{y,k}^T}{g_{y,k}^T s_{y,k}}) + \frac{s_{y,k} s_{y,k}^T}{g_{y,k}^T s_{y,k}}$
    \end{algorithmic}
\end{algorithm}

\section{Experiment}
\begin{itemize}
    \item In gaussian mean, cqn perform better than gda, worse than cn, but gdqn and qngd perform similar and worse than gda
    \item In gaussian mean ill, gdqn and cqn perform better than gda but worse than cq, qngd perform worse than gda
    \item In gaussian covariance, diverge
\end{itemize}

\section{Other Note}
\begin{itemize}
    \item generator is the minimizer which is the slow leader
    \item discriminator is the maximizer which is the faster follower
    \item will the update of other side effect the approximation of Hessian
    \item run the inner problem more times and the Hessian should not be used between iterations
    \item care about the input of inner problem: - loss fn
    \item bug: rho when calculating V and H
    \item take notes
\end{itemize}

\end{document}